\paragraph{An operator construction of the common double cover.}
Here is a direct justification that the same lifts may be used in
the comparison. It also fixes the signs of the diagonal operators.
Work over $k=\mathbb Q_p$, including $p=2$, or $k=\mathbb R$,
with the additive character just specified. For a nondegenerate
ternary $Q$ write $\gamma_Q=\gamma(Q)$ and define
\[
 \begin{split}
 U_Q(b)\phi(x)&=\mathbf e_k(bQ(x))\phi(x),\\
 D(a)\phi(x)&=|a|_k^{3/2}\phi(ax),\\
 W_Q&=\gamma_Q\mathcal F_Q,\qquad
 M_Q(a)=r_Q(a)D(a),\qquad
 r_Q(a)=\frac{\gamma(aQ)}{\gamma(Q)}.
 \end{split}
\]
The Weyl matrix in this convention is
$w=\left(\begin{smallmatrix}0&1\\-1&0\end{smallmatrix}\right)$.
The opposite Fourier-sign convention changes both compared
representations together.

The Fourier formula of \cite[I.14, Theorem~2]{weil}, or its
completion-of-the-square calculation, says for $b\ne0$ that
\[
 \int_{k^3}
 \mathbf e_k\bigl(bQ(z)+B_Q(z,v)\bigr)\,dz_Q
   =|b|_k^{-3/2}\gamma(bQ)\mathbf e_k(-b^{-1}Q(v)).
\]
The integral has the oscillatory meaning of that theorem.
Applying it to a Schwartz function gives the operator identity
\[
 \mathcal F_Q U_Q(b)\mathcal F_Q
 =\gamma(bQ)U_Q(-b^{-1})D(b^{-1})
                   \mathcal F_Q^{-1}U_Q(-b^{-1}).
\]
Thus no choice of an individual square root of a Gauss sum is
needed. The index identities proved above also give
\[
 \gamma(Q)^2=(\det Q,-1)_k\gamma(1)^6,
 \qquad \gamma(1)^4=(-1,-1)_k,
 \qquad r_Q(a)r_Q(b)=(a,b)_k r_Q(ab).
\]
For the middle equality apply the quaternion-norm formula to
$\langle1,1,1,1\rangle$; the first follows by applying the same
formula to $\langle1,-a_i,1,-a_i\rangle$ for each diagonal
entry $a_i$ of $Q$.

Consequently the displayed operators satisfy
\[
 \begin{gathered}
 U_Q(b)U_Q(c)=U_Q(b+c),\qquad
 M_Q(a)U_Q(b)M_Q(a)^{-1}=U_Q(a^2b),\\
 M_Q(a)M_Q(b)=(a,b)_kM_Q(ab),\qquad
 W_QM_Q(a)W_Q^{-1}=M_Q(a^{-1}),\\
 W_Q^2=(-1,-1)_kM_Q(-1),\\
 W_QU_Q(b)W_Q
 =(-b,-1)_kU_Q(-b^{-1})M_Q(-b^{-1})W_QU_Q(-b^{-1})
 \quad(b\ne0).
 \end{gathered}
\]
For the last equality, the preceding Fourier identity leaves
the scalar $\gamma(Q)^2\gamma(bQ)^2$; the index-square formula
reduces it to $(-b,-1)_k$. Every scalar in these relations is
independent of the ternary form. Interpreting a scalar $-1$ as
the central element $\varepsilon$ gives one common double
cover of $\operatorname{SL}_2(k)$. Indeed these relations reduce
every word to one of
$\varepsilon^eU(x)M(a)$ or
$\varepsilon^eU(x)M(a)WU(y)$, with $e\in\{0,1\}$.
The ordinary Bruhat decomposition makes their projections
unique, apart from $\varepsilon^e$; the operators send
$\varepsilon$ to $-\operatorname{id}$, so the central kernel
has exactly two elements. This realizes the rank-three pullback
of Weil's metaplectic operators, rather than assuming a section
over the diagonal units.

For $Q_0=\langle1,1,1\rangle$ these operators are the tensor
products of the three ordinary one-dimensional theta operators.
For either $Q$ or $Q_0$, changing a lift multiplies its operator
by the same $-1$. Hence the quotient of their eigenvalues is
independent of the lift. At two the relation
$M_Q(-1)^2=(-1,-1)_2\operatorname{id}=-\operatorname{id}$
explicitly retains the order-four lift of $-I$.

On a unit $a$ the quotient of the two diagonal index factors is
\[
 \frac{r_Q(a)}{r_{Q_0}(a)}=(a,\det Q)_k.
\]
Together with $f_p(ax)=\alpha_p(a)f_p(x)$, this is exactly
the eigenvalue quotient $\alpha_p(a)(a,\det Q)_p$ used below.
There is no remaining eighth-root factor at two.

